\documentclass[reqno]{exam}

\usepackage{amssymb, amsmath, amsthm, stmaryrd}
\usepackage{fullpage, parskip}
\usepackage{enumitem}
\usepackage{tikz}
\usepackage{ulem}
\usepackage{hyperref}

\title{Introduction to Computational Math Spring 2026 \\ Homework 02}
\date{Due: Friday, Feb 06, 2025, 11:59 PM}
\author{}

\begin{document}
\maketitle
\thispagestyle{empty}

Textbook = Driscoll and Braun (\url{https://fncbook.com/})

\begin{enumerate}

\item Textbook Exercise 1.2.3

\item Textbook Exercise 1.2.6

\item Textbook Exercise 1.2.8

\item Textbook Exercise 1.4.1

\item

Jupyter notebook.

\end{enumerate}

See instructions on the next page for details on submission and grading policy.

\newpage
\section*{Homework Grading Policy}

Your homework will be graded based on ``honest attempts''.
This means that you will not lose points for mistakes as long as you have attempted each problem with sufficient detail and effort.
The purpose of this policy is to:
\begin{enumerate}
  \item Encourage you to solve problems independently without relying on AI tools for the mathematical work.
  \item Avoid penalizing students who work without artificial assistance and may make more mistakes as part of the learning process.
\end{enumerate}

\section*{Quiz Information}

The quiz will be held on {\bf Tuesday, February 10, 2026} during the discussion section.

Quiz questions will be modifications of the homework problems and material discussed in class. Students who rely on automated solutions to complete this homework will likely struggle on the quiz, as it requires genuine understanding of the concepts.

\textbf{Topics covered:} 

Chapters 1.2, 1.4. The emphasis will be on Ch 1.2.


\section*{Submission Instructions}

\begin{itemize}
  \item All homework (written and coding) must be submitted on Gradescope as a \textbf{single PDF} file.

  \item For the Jupyter Notebook component:
        \begin{itemize}
          \item Follow all instructions in the notebook.
          \item Include relevant lines of code, outputs, and any requested plots.
          \item Respond to questions using markdown cells where applicable.
          \item \textbf{Note:} You are welcome to use AI tools (such as HopGPT or Copilot) for programming tasks, as coding proficiency is not being tested in this course. However, ensure you understand the outputs and can interpret the results.
        \end{itemize}

  \item Converting your notebook to PDF:
        \begin{itemize}
          \item \textbf{Jupyter:} Open your notebook in a web browser and use the browser's ``Print to PDF'' function.

          \item \textbf{VS Code:} Export your finished notebook to HTML, open it in a browser, and use ``Print to PDF.''

          \item \textbf{Quarto:} Use the command \texttt{quarto render notebook.ipynb --to pdf} to directly convert your notebook to PDF.
        \end{itemize}

  \item Ensure your PDF is properly formatted and not excessively tall or narrow. Check that all content is clearly visible before submitting.
\end{itemize}
\end{document}